% --- Archon physics formalization source begin ---
% archon:physics
% archon:covers PhyXMiniProblems/problem_phyx_mini_0054.lean
% archon:source-report reports/phyx_mini/problem_phyx_mini_0054.source.json

\chapter{Physics problem phyx\_mini\_0054}
\label{ch:PhyXMiniProblems_problem_phyx_mini_0054}

\paragraph{Problem source.}
\#\# Physical scenario

One end of a 4.0-cm-diameter glass rod is shaped like a hemisphere. A small lightbulb is 6.0 cm from the end of the rod.

\#\# Question

How faraway inside is the real image located in the rod?

\#\# Answer choices

- A: A: 14cm

- B: B: 18cm

- C: C: 96cm

- D: D: 74cm

\#\# Figure caption (auxiliary; use the image as primary evidence)

The image shows a diagram illustrating the refraction of light through a cylindrical lens. Here's a detailed description:

- **Objects and Scenes**:
  - There is a black dot labeled "Object" on the left side of the image.
  - Multiple rays are drawn from the object, converging through a shaded lens.
  - The lens is cylindrical in shape, with shading indicating its presence.

- **Relationships**:
  - The rays travel from the object, refract at the lens surface, and converge on the right side.
  - The point where the rays meet is labeled "Image" with a black dot.

- **Text and Labels**:
  - "n₁ = 1.00" is written above the object, indicating the refractive index of the medium outside the lens.
  - "n₂ = 1.50" is written above the lens, representing its refractive index.
  - The text "R = 40 cm" is annotated next to a double-headed arrow indicating the radius of curvature of the lens.
  - "s = 60 cm" is labeled beneath the object, indicating the object distance from the lens.
  - "s'" is labeled below the arrow pointing to the image, indicating the image distance from the lens.

This diagram is likely used to explain the principles of refraction and image formation through lenses.

\#\# Dataset metadata

Geometrical Optics; Physical Model Grounding Reasoning, Multi-Formula Reasoning

\paragraph{Recorded answer/context.}
B: B: 18cm

\paragraph{Figure/image path.}
/root/proposal\_for\_physic/science-mango/phyx\_mini\_run/phyx\_data/test\_image/54.png

\paragraph{Formalization target.}
create a compiling Lean file with sorry bodies at `PhyXMiniProblems/problem\_phyx\_mini\_0054.lean`. The Lean declarations must preserve the physical quantities, dimensions or dimensional roles, figure labels, governing-law hypotheses, and final relation expressed by this problem.
Use Mathlib/Physlib names found through LeanExplore where available. If a physics API is missing, introduce faithful local abstractions rather than scalar placeholder aliases.

\begin{theorem}[Physics formalization target]
\label{thm:physics:phyx_mini_0054:target}
\lean{PhyXMiniProblems.ProblemPhyXMini0054.problem_phyx_mini_0054}
\uses{def:physics:phyx-mini-0054:phyxminiproblems-problemphyxmini0054-hemisphericalglassrodsetup, def:physics:phyx-mini-0054:phyxminiproblems-problemphyxmini0054-imagedistanceincentimeters, def:physics:phyx-mini-0054:phyxminiproblems-problemphyxmini0054-matchesproblemandprimaryfigure, def:physics:phyx-mini-0054:phyxminiproblems-problemphyxmini0054-modelshemisphericalrodend, def:physics:phyx-mini-0054:phyxminiproblems-problemphyxmini0054-hasdepictedphysicalconfiguration, def:physics:phyx-mini-0054:phyxminiproblems-problemphyxmini0054-obeysparaxialsphericalsurfaceequation, def:physics:phyx-mini-0054:phyxminiproblems-problemphyxmini0054-matchesanswerchoice}
The assigned autoformalize agent should translate this physics problem into Lean declarations in the covered file, with theorem and lemma proof bodies written as `by sorry`.
\end{theorem}
\begin{proof}
This is an autoformalization task, not a proof task. Produce statements that can later be proved by the physics prover without weakening the physical model.
\end{proof}

% --- Archon declaration topology begin ---
\section{Declaration topology}
\begin{definition}[Length Quantity]
  \label{def:physics:phyx-mini-0054:phyxminiproblems-problemphyxmini0054-lengthquantity}
  \lean{PhyXMiniProblems.ProblemPhyXMini0054.LengthQuantity}
  A signed physical length whose readout changes coherently with unit choice
\end{definition}
\begin{proof}
  This is the declared model definition of Length Quantity.
\end{proof}
\begin{definition}[centimeter Unit Choices]
  \label{def:physics:phyx-mini-0054:phyxminiproblems-problemphyxmini0054-centimeterunitchoices}
  \lean{PhyXMiniProblems.ProblemPhyXMini0054.centimeterUnitChoices}
  SI base units with centimeters selected as the length unit
\end{definition}
\begin{proof}
  This is the declared model definition of centimeter Unit Choices.
\end{proof}
\begin{definition}[length In Centimeters]
  \label{def:physics:phyx-mini-0054:phyxminiproblems-problemphyxmini0054-lengthincentimeters}
  \lean{PhyXMiniProblems.ProblemPhyXMini0054.lengthInCentimeters}
  \uses{def:physics:phyx-mini-0054:phyxminiproblems-problemphyxmini0054-lengthquantity, def:physics:phyx-mini-0054:phyxminiproblems-problemphyxmini0054-centimeterunitchoices}
  The signed scalar readout of a physical length in centimeters
\end{definition}
\begin{proof}
  This is the declared model definition of length In Centimeters.
\end{proof}
\begin{definition}[Optical Region]
  \label{def:physics:phyx-mini-0054:phyxminiproblems-problemphyxmini0054-opticalregion}
  \lean{PhyXMiniProblems.ProblemPhyXMini0054.OpticalRegion}
  The two optical regions separated by the hemispherical rod surface
\end{definition}
\begin{proof}
  This is the declared model definition of Optical Region.
\end{proof}
\begin{definition}[Figure Point]
  \label{def:physics:phyx-mini-0054:phyxminiproblems-problemphyxmini0054-figurepoint}
  \lean{PhyXMiniProblems.ProblemPhyXMini0054.FigurePoint}
  The three labeled points on the optical axis in the primary figure
\end{definition}
\begin{proof}
  This is the declared model definition of Figure Point.
\end{proof}
\begin{definition}[Geometrical Image Kind]
  \label{def:physics:phyx-mini-0054:phyxminiproblems-problemphyxmini0054-geometricalimagekind}
  \lean{PhyXMiniProblems.ProblemPhyXMini0054.GeometricalImageKind}
  Whether the depicted image is formed by rays or by backward extensions
\end{definition}
\begin{proof}
  This is the declared model definition of Geometrical Image Kind.
\end{proof}
\begin{definition}[Hemispherical Glass Rod Setup]
  \label{def:physics:phyx-mini-0054:phyxminiproblems-problemphyxmini0054-hemisphericalglassrodsetup}
  \lean{PhyXMiniProblems.ProblemPhyXMini0054.HemisphericalGlassRodSetup}
  \uses{def:physics:phyx-mini-0054:phyxminiproblems-problemphyxmini0054-lengthquantity, def:physics:phyx-mini-0054:phyxminiproblems-problemphyxmini0054-opticalregion, def:physics:phyx-mini-0054:phyxminiproblems-problemphyxmini0054-figurepoint, def:physics:phyx-mini-0054:phyxminiproblems-problemphyxmini0054-geometricalimagekind}
  The dimensionful and dimensionless quantities in the rod setup. The surface radius is positive for the convex surface whose center of curvature lies inside the glass. The axial positions increase from the bulb, through the surface vertex, into the rod
\end{definition}
\begin{proof}
  This is the declared model definition of Hemispherical Glass Rod Setup.
\end{proof}
\begin{definition}[axial Separation Readout]
  \label{def:physics:phyx-mini-0054:phyxminiproblems-problemphyxmini0054-axialseparationreadout}
  \lean{PhyXMiniProblems.ProblemPhyXMini0054.axialSeparationReadout}
  \uses{def:physics:phyx-mini-0054:phyxminiproblems-problemphyxmini0054-figurepoint, def:physics:phyx-mini-0054:phyxminiproblems-problemphyxmini0054-hemisphericalglassrodsetup}
  Signed separation of two figure points, read in a specified unit system
\end{definition}
\begin{proof}
  This is the declared model definition of axial Separation Readout.
\end{proof}
\begin{definition}[object Distance Readout]
  \label{def:physics:phyx-mini-0054:phyxminiproblems-problemphyxmini0054-objectdistancereadout}
  \lean{PhyXMiniProblems.ProblemPhyXMini0054.objectDistanceReadout}
  \uses{def:physics:phyx-mini-0054:phyxminiproblems-problemphyxmini0054-hemisphericalglassrodsetup, def:physics:phyx-mini-0054:phyxminiproblems-problemphyxmini0054-axialseparationreadout}
  The positive object-distance readout s, from the bulb to the vertex
\end{definition}
\begin{proof}
  This is the declared model definition of object Distance Readout.
\end{proof}
\begin{definition}[image Distance Readout]
  \label{def:physics:phyx-mini-0054:phyxminiproblems-problemphyxmini0054-imagedistancereadout}
  \lean{PhyXMiniProblems.ProblemPhyXMini0054.imageDistanceReadout}
  \uses{def:physics:phyx-mini-0054:phyxminiproblems-problemphyxmini0054-hemisphericalglassrodsetup, def:physics:phyx-mini-0054:phyxminiproblems-problemphyxmini0054-axialseparationreadout}
  The positive real-image-distance readout s', from the vertex into the rod
\end{definition}
\begin{proof}
  This is the declared model definition of image Distance Readout.
\end{proof}
\begin{definition}[object Distance In Centimeters]
  \label{def:physics:phyx-mini-0054:phyxminiproblems-problemphyxmini0054-objectdistanceincentimeters}
  \lean{PhyXMiniProblems.ProblemPhyXMini0054.objectDistanceInCentimeters}
  \uses{def:physics:phyx-mini-0054:phyxminiproblems-problemphyxmini0054-centimeterunitchoices, def:physics:phyx-mini-0054:phyxminiproblems-problemphyxmini0054-hemisphericalglassrodsetup, def:physics:phyx-mini-0054:phyxminiproblems-problemphyxmini0054-objectdistancereadout}
  The object distance in the centimeter unit system used by the figure
\end{definition}
\begin{proof}
  This is the declared model definition of object Distance In Centimeters.
\end{proof}
\begin{definition}[image Distance In Centimeters]
  \label{def:physics:phyx-mini-0054:phyxminiproblems-problemphyxmini0054-imagedistanceincentimeters}
  \lean{PhyXMiniProblems.ProblemPhyXMini0054.imageDistanceInCentimeters}
  \uses{def:physics:phyx-mini-0054:phyxminiproblems-problemphyxmini0054-centimeterunitchoices, def:physics:phyx-mini-0054:phyxminiproblems-problemphyxmini0054-hemisphericalglassrodsetup, def:physics:phyx-mini-0054:phyxminiproblems-problemphyxmini0054-imagedistancereadout}
  The requested real-image distance in centimeters
\end{definition}
\begin{proof}
  This is the declared model definition of image Distance In Centimeters.
\end{proof}
\begin{definition}[Matches Problem And Primary Figure]
  \label{def:physics:phyx-mini-0054:phyxminiproblems-problemphyxmini0054-matchesproblemandprimaryfigure}
  \lean{PhyXMiniProblems.ProblemPhyXMini0054.MatchesProblemAndPrimaryFigure}
  \uses{def:physics:phyx-mini-0054:phyxminiproblems-problemphyxmini0054-lengthincentimeters, def:physics:phyx-mini-0054:phyxminiproblems-problemphyxmini0054-hemisphericalglassrodsetup, def:physics:phyx-mini-0054:phyxminiproblems-problemphyxmini0054-objectdistanceincentimeters}
  Problem-statement and primary-figure data. The image is explicitly depicted as a real convergence point in the glass. No numerical value for its distance is included here
\end{definition}
\begin{proof}
  This is the declared model definition of Matches Problem And Primary Figure.
\end{proof}
\begin{definition}[Models Hemispherical Rod End]
  \label{def:physics:phyx-mini-0054:phyxminiproblems-problemphyxmini0054-modelshemisphericalrodend}
  \lean{PhyXMiniProblems.ProblemPhyXMini0054.ModelsHemisphericalRodEnd}
  \uses{def:physics:phyx-mini-0054:phyxminiproblems-problemphyxmini0054-hemisphericalglassrodsetup}
  Geometric meaning of “one end of the rod is shaped like a hemisphere”: the spherical surface radius is half the rod diameter. This is setup geometry, not the requested image-distance result
\end{definition}
\begin{proof}
  This is the declared model definition of Models Hemispherical Rod End.
\end{proof}
\begin{definition}[Has Depicted Physical Configuration]
  \label{def:physics:phyx-mini-0054:phyxminiproblems-problemphyxmini0054-hasdepictedphysicalconfiguration}
  \lean{PhyXMiniProblems.ProblemPhyXMini0054.HasDepictedPhysicalConfiguration}
  \uses{def:physics:phyx-mini-0054:phyxminiproblems-problemphyxmini0054-lengthincentimeters, def:physics:phyx-mini-0054:phyxminiproblems-problemphyxmini0054-opticalregion, def:physics:phyx-mini-0054:phyxminiproblems-problemphyxmini0054-hemisphericalglassrodsetup, def:physics:phyx-mini-0054:phyxminiproblems-problemphyxmini0054-objectdistanceincentimeters, def:physics:phyx-mini-0054:phyxminiproblems-problemphyxmini0054-imagedistanceincentimeters}
  Positivity and axial-order conditions for the physical branch in the figure
\end{definition}
\begin{proof}
  This is the declared model definition of Has Depicted Physical Configuration.
\end{proof}
\begin{definition}[Obeys Paraxial Spherical Surface Equation]
  \label{def:physics:phyx-mini-0054:phyxminiproblems-problemphyxmini0054-obeysparaxialsphericalsurfaceequation}
  \lean{PhyXMiniProblems.ProblemPhyXMini0054.ObeysParaxialSphericalSurfaceEquation}
  \uses{def:physics:phyx-mini-0054:phyxminiproblems-problemphyxmini0054-hemisphericalglassrodsetup, def:physics:phyx-mini-0054:phyxminiproblems-problemphyxmini0054-objectdistancereadout, def:physics:phyx-mini-0054:phyxminiproblems-problemphyxmini0054-imagedistancereadout}
  The Gaussian paraxial equation for refraction at one convex spherical surface, written with positive object distance s, positive real-image distance s', and positive radius R: n₁ / s + n₂ / s' = (n₂ - n₁) / R. It is required for every unit choice, so it is a dimensionally coherent governing law rather than a formula specialized to the requested answer
\end{definition}
\begin{proof}
  This is the declared model definition of Obeys Paraxial Spherical Surface Equation.
\end{proof}
\begin{definition}[Answer Choice]
  \label{def:physics:phyx-mini-0054:phyxminiproblems-problemphyxmini0054-answerchoice}
  \lean{PhyXMiniProblems.ProblemPhyXMini0054.AnswerChoice}
  Labels of the four displayed answer choices
\end{definition}
\begin{proof}
  This is the declared model definition of Answer Choice.
\end{proof}
\begin{definition}[answer Image Distance In Centimeters]
  \label{def:physics:phyx-mini-0054:phyxminiproblems-problemphyxmini0054-answerimagedistanceincentimeters}
  \lean{PhyXMiniProblems.ProblemPhyXMini0054.answerImageDistanceInCentimeters}
  \uses{def:physics:phyx-mini-0054:phyxminiproblems-problemphyxmini0054-answerchoice}
  The centimeter readout printed beside each answer choice
\end{definition}
\begin{proof}
  This is the declared model definition of answer Image Distance In Centimeters.
\end{proof}
\begin{definition}[Matches Answer Choice]
  \label{def:physics:phyx-mini-0054:phyxminiproblems-problemphyxmini0054-matchesanswerchoice}
  \lean{PhyXMiniProblems.ProblemPhyXMini0054.MatchesAnswerChoice}
  \uses{def:physics:phyx-mini-0054:phyxminiproblems-problemphyxmini0054-hemisphericalglassrodsetup, def:physics:phyx-mini-0054:phyxminiproblems-problemphyxmini0054-imagedistanceincentimeters, def:physics:phyx-mini-0054:phyxminiproblems-problemphyxmini0054-answerchoice, def:physics:phyx-mini-0054:phyxminiproblems-problemphyxmini0054-answerimagedistanceincentimeters}
  Agreement of the physical image-distance readout with a displayed choice
\end{definition}
\begin{proof}
  This is the declared model definition of Matches Answer Choice.
\end{proof}
\begin{lemma}[surface Radius In Centimeters eq two]
  \label{lem:physics:phyx-mini-0054:phyxminiproblems-problemphyxmini0054-surfaceradiusincentimeters-eq-two}
  \lean{PhyXMiniProblems.ProblemPhyXMini0054.surfaceRadiusInCentimeters_eq_two}
  \uses{def:physics:phyx-mini-0054:phyxminiproblems-problemphyxmini0054-lengthincentimeters, def:physics:phyx-mini-0054:phyxminiproblems-problemphyxmini0054-hemisphericalglassrodsetup, def:physics:phyx-mini-0054:phyxminiproblems-problemphyxmini0054-matchesproblemandprimaryfigure, def:physics:phyx-mini-0054:phyxminiproblems-problemphyxmini0054-modelshemisphericalrodend}
  The 4.0 cm diameter and hemispherical geometry give R = 2.0 cm
\end{lemma}
\begin{proof}
  The conclusion follows from the governing relations and figure readouts named in the statement of surface Radius In Centimeters eq two.
\end{proof}
% --- Archon declaration topology end ---
% --- Archon physics formalization source end ---
